\documentclass[11pt]{article}
\usepackage[a4paper,margin=1in]{geometry}
\usepackage[utf8]{inputenc}
\usepackage[T1]{fontenc}
\usepackage{amsmath,amssymb}
\usepackage{booktabs}
\usepackage{xcolor}
\usepackage{hyperref}

\newcommand{\pending}[1]{\textcolor{red}{\textbf{[PENDING: #1]}}}

\title{\textbf{What a Benchmark Can Resolve: Training Noise as a Ceiling on\\
Architecture Comparison and on Predictor Evaluation}}
\author{Vasili Gavrilov\thanks{Independent researcher. Code and derived tables:
\url{https://github.com/Anode1/SMBPANN}.}}
\date{August 2026}

\begin{document}
\maketitle

\begin{abstract}
We measure what a neural architecture benchmark can and cannot resolve, using the repeated training runs
its authors performed and its distributed tables average away. NAS-Bench-101 trained each of its 423,624
architectures three times; NAS-Bench-201 trained each of 15,625 architectures three times on each of
three datasets, at every epoch up to 200. Recovering those runs gives a reference that is statistically
independent of any single one of them, and four things follow.

First, comparing two architectures from one training run each disagrees with an independent reference
33 to 42 per cent of the time when the observed difference is 0.1 to 0.2 percentage points, on both
benchmarks and all three datasets. Second, and worse for practice, at low training budget the comparison
is a coin flip at every effect size: at 4 epochs on NAS-Bench-101 a one-to-two point difference is called
backwards 45 per cent of the time, against 4 per cent at 108 epochs, which bears directly on
multi-fidelity and early-stopping search. Third, a benchmark's reported optimum is not a point: 3,558 of
NAS-Bench-101's architectures cannot be separated from its winner using its own data, so regret measured
against that winner inherits the width. Fourth, and most useful to the largest audience, there is a
ceiling on the rank correlation any performance predictor can achieve, because the label it is scored
against is itself a noisy average. Among the top 1,000 architectures of NAS-Bench-101 that ceiling is
$\tau = 0.084$; the same quantity on NAS-Bench-201 is $0.536$.

We then show these constants are largely not properties of the individual benchmarks. Writing $\sigma_W$
for training noise and $\sigma_B$ for the spread of true quality over a set of architectures, the ceiling
follows from the single ratio $\sigma_W/\sigma_B$, and the collapse of the ceiling on high-performing
subsets follows from range restriction shrinking $\sigma_B$ while $\sigma_W$ stays fixed. Across 92 cells
spanning two benchmarks, three datasets, ten training budgets and four subset sizes, a one-parameter
prediction matches the measured ceiling to within $0.07$ while $\sigma_W/\sigma_B < 0.6$, degrading to
about $0.15$ beyond it. The single large failure is diagnosable and instructive: NAS-Bench-101's noise is
a mixture, since roughly one architecture in a hundred sometimes trains and sometimes collapses to chance
accuracy, and excluding those moves the prediction error from $+0.41$ to $+0.02$. The practical
consequence is that the ceiling can be estimated for any benchmark or evaluation set from two variances,
without the replicates at scale that almost nobody has.
\end{abstract}

\section{Why this is measurable at all}\label{sec:why}

Training a network is random. Starting weights are drawn, examples are shuffled, and on the usual
hardware the reductions are not associative, so the same architecture trained twice gives two different
accuracies. Any single accuracy is therefore the architecture's quality plus a draw of noise. This is not
controversial and it is not new; what is missing is a number.

The number is hard to obtain because it requires training the same thing repeatedly, which nobody can
afford at modern scale. Two public benchmarks, however, already paid that cost for a different reason.
NAS-Bench-101 trained each of 423,624 cell architectures three separate times on CIFAR-10, at each of
four epoch budgets. NAS-Bench-201 trained each of 15,625 architectures three times on each of CIFAR-10,
CIFAR-100 and a downscaled ImageNet, recording every epoch to 200. Both are among the very few public
datasets containing genuine training replicates at scale.

Both distributed tables average those runs into one number per architecture. That is the right table for
ranking architectures and it destroys the only structure a study of the noise can use, so the first step
here is recovering the individual runs. We do so and check the recovery: the mean of our three recovered
runs reproduces the distributed averaged table to within $10^{-4}$ percentage points on all 423,624 rows
of NAS-Bench-101.

With three runs in hand, every measurement below uses the same device. One run supplies the judgment; the
remaining runs, averaged, supply the reference. The two use disjoint runs, so the reference is
statistically independent of the judgment, and no quantity is compared against itself.

\section{Theory: one ratio should govern the ceiling}\label{sec:theory}

Let architecture $i$ have true quality $q_i$, with variance $\sigma_B^2$ over whatever set of
architectures is under consideration, and let a training run return $x_{ij} = q_i + e_{ij}$ with
$\mathrm{Var}(e) = \sigma_W^2$. A one-run label $A_i = x_{i1}$ has variance $\sigma_B^2 + \sigma_W^2$; a
two-run label $B_i = (x_{i2}+x_{i3})/2$ has variance $\sigma_B^2 + \sigma_W^2/2$; and since the noise
terms are independent, $\mathrm{Cov}(A_i,B_i) = \sigma_B^2$. Hence
\begin{equation}
  \rho \;=\; \frac{\sigma_B^2}{\sqrt{(\sigma_B^2+\sigma_W^2)\,(\sigma_B^2+\sigma_W^2/2)}}
        \;=\; \frac{1}{\sqrt{(1+r^2)(1+r^2/2)}}, \qquad r \equiv \sigma_W/\sigma_B ,
  \label{eq:rho}
\end{equation}
and for jointly normal variables Kendall's $\tau = (2/\pi)\arcsin\rho$. Two consequences are worth stating
before any data is seen.

The ceiling on rank correlation is a function of the single ratio $r$, not an independent property of a
benchmark. And the ceiling must fall on high-performing subsets, because selecting the top architectures
restricts the range of $q$ and so shrinks $\sigma_B$, while $\sigma_W$ is unaffected by which
architectures we chose to look at. The severity of that collapse is then predicted rather than merely
observed.

The same two quantities give a resolution: the standard error of a difference between two architectures
each trained $n$ times is $\sigma_W\sqrt{2/n}$, so a true gap $\Delta$ is resolvable at 95 per cent
confidence when $\Delta > 1.96\,\sigma_W\sqrt{2/n}$, and at $n=1$ the floor is $2.77\,\sigma_W$.

\section{What we measure}\label{sec:results}

\subsection{A single run disagrees with an independent reference}

Binning by the \emph{observed} single-run difference, which is what a reader of a paper has in front of
them, over two million random pairs per dataset:

\begin{center}\small
\begin{tabular}{lcccc}
\toprule
observed difference & NB101 C10 & NB201 C10 & NB201 C100 & NB201 IN16 \\
\midrule
$0.05$--$0.10$ pp & $46.1\%$ & $39.7\%$ & $42.3\%$ & $42.3\%$ \\
$0.10$--$0.20$     & $41.6$   & $33.0$   & $38.2$   & $38.4$ \\
$0.20$--$0.30$     & $36.4$   & $25.3$   & $33.4$   & $34.2$ \\
$0.50$--$1.00$     & $16.3$   & $4.3$    & $14.5$   & $15.9$ \\
$1.00$--$2.00$     & $4.4$    & $0.3$    & $3.2$    & $4.2$ \\
\bottomrule
\end{tabular}
\end{center}

The rates track the noise level across datasets, which is the ordering a correct measurement should
produce. Two cautions. With three runs there is no noiseless reference, so this is a replication
disagreement rate rather than a decomposition isolating one run's own error. And because the bins are
defined by the observed difference, each rate is a predictive posterior, ``I measured $0.15$, how much
should I believe it'', and not a statement that a given fraction of true $0.15$ gaps are called wrongly.

\subsection{At low training budget the comparison carries no information}

NAS-Bench-101's full archive holds three runs at each of 4, 12, 36 and 108 epochs.

\begin{center}\small
\begin{tabular}{lccccc}
\toprule
epochs & median $\sigma_W$ & $0.1$--$0.2$ & $0.2$--$0.3$ & $0.5$--$1.0$ & $1.0$--$2.0$ \\
\midrule
4   & $4.147$ & $48.9\%$ & $51.2\%$ & $47.0\%$ & $45.1\%$ \\
12  & $5.090$ & $49.2$   & $46.7$   & $47.3$   & $43.5$ \\
36  & $0.804$ & $47.0$   & $44.5$   & $30.8$   & $18.3$ \\
108 & $0.328$ & $41.6$   & $36.3$   & $16.5$   & $4.4$ \\
\bottomrule
\end{tabular}
\end{center}

At 4 and 12 epochs the comparison is a coin flip at every effect size up to two percentage points: a
gap that is reliable after full training, $4.4$ per cent backwards, is $45$ per cent backwards at 4
epochs. This bears directly on multi-fidelity and early-stopping search, which is standard practice and
which ranks candidates on exactly these cheap evaluations on the assumption that the ranking transfers.
Note also that the noise is not monotone in budget, 12 epochs being noisier than 4, plausibly because
4-epoch accuracy is uniformly poor and leaves less room to diverge.

\subsection{The reported optimum is a set, not a point}

Testing each architecture against the benchmark's best using the pair's own noise:

\begin{center}\small
\begin{tabular}{lcc}
\toprule
benchmark / dataset & best 3-run mean & not separable at 95\% \\
\midrule
NAS-Bench-101 CIFAR-10        & $95.06$ & $3{,}558$ of $423{,}624$ \\
NAS-Bench-201 CIFAR-10        & $94.37$ & $34$ of $15{,}625$ \\
NAS-Bench-201 CIFAR-100       & $73.50$ & $4$ of $15{,}625$ \\
NAS-Bench-201 ImageNet16-120  & $46.88$ & $15$ of $15{,}625$ \\
\bottomrule
\end{tabular}
\end{center}

Regret measured against a single reported optimum inherits that width. NAS-Bench-201 is far better
behaved on this axis than NAS-Bench-101, which is worth stating separately rather than averaging over.

\subsection{The ceiling on predictor evaluation}

Zero-cost proxies, surrogates and performance predictors are all scored by rank correlation against a
benchmark label that is the mean of a few noisy runs, so no predictor can exceed the correlation the
label achieves with itself. Measured as $\tau$ between a one-run label and an independent two-run label,
with the subset chosen by the reference alone so that the judgment run plays no part in the selection:

\begin{center}\small
\begin{tabular}{lcccc}
\toprule
subset & NB101 C10 & NB201 C10 & NB201 C100 & NB201 IN16 \\
\midrule
top 100    & $0.142$ & $0.384$ & $0.596$ & $0.419$ \\
top 1000   & $0.084$ & $0.536$ & $0.549$ & $0.495$ \\
top 10000  & $0.186$ & $0.849$ & $0.845$ & $0.889$ \\
all        & $0.832$ & $0.923$ & $0.924$ & $0.941$ \\
\bottomrule
\end{tabular}
\end{center}

The whole-set values all exceed $0.83$, which is why the problem is easy to miss: over the full space a
single run ranks architectures well, because most of the space is separated by margins far larger than
the noise. The top-$k$ rows are the relevant ones, since a predictor's job is to rank good architectures
rather than to separate good from broken, and there the ceiling on NAS-Bench-101 is $0.08$ to $0.19$. A
published correlation near or above that, computed on a comparable subset, is not evidence of a good
predictor.

\section{The constants are mostly consequences of one ratio}\label{sec:law}

Equation~\eqref{eq:rho} predicts each of those ceilings from $\sigma_W/\sigma_B$ measured on the same
subset. We estimate $\sigma_W$ as the pooled within-architecture value and $\sigma_B$ from the judgment
run alone, whose variance is $\sigma_B^2+\sigma_W^2$; using the three-run mean here would be wrong,
because the subset was chosen using the reference and the mean contains it.

Across 92 cells, spanning two benchmarks, three datasets, ten training budgets and four subset sizes,
with $\tau$ from $0.42$ to $0.94$ and $r$ from $0.03$ to $2.6$:

\begin{center}\small
\begin{tabular}{lccc}
\toprule
$\sigma_W/\sigma_B$ & cells & mean residual & RMS residual \\
\midrule
$<0.1$      & 3  & $-0.038$ & $0.042$ \\
$0.1$--$0.3$ & 16 & $-0.049$ & $0.064$ \\
$0.3$--$0.6$ & 20 & $-0.035$ & $0.075$ \\
$0.6$--$1.0$ & 27 & $-0.074$ & $0.132$ \\
$>1.0$      & 26 & $-0.129$ & $0.160$ \\
\bottomrule
\end{tabular}
\end{center}

So a one-parameter prediction reproduces the ceiling to within about $0.07$ while the noise is below
sixty per cent of the quality spread, and degrades beyond that with a systematic tendency to
over-predict.

\textbf{The single large failure is the informative one.} On the full NAS-Bench-101 set the prediction is
$0.42$ against a measured $0.832$, a residual of $+0.41$. The cause is that its noise is not one
distribution but a mixture: roughly one architecture in a hundred sometimes trains and sometimes
collapses to chance accuracy, with $\sigma_W$ up to $48$ percentage points. Those architectures dominate
the pooled $\sigma_W$ without dominating the rank correlation, because the well-behaved remainder still
orders consistently. Excluding them moves $r$ from $0.924$ to $0.246$ and the residual from $+0.41$ to
$+0.02$. The residual is therefore not merely an error term; it is a detector for heterogeneous training
noise, and its sign says which way the heterogeneity runs.

\textbf{Why this matters more than the constants.} If the ceiling followed from each benchmark's
idiosyncrasies, it would have to be measured wherever it was wanted, which requires replicates at scale
that almost nobody has. Because it follows from two variances, it can be estimated from a handful of
repeated evaluations, or in some settings from a variance decomposition of a single run's own
sub-samples. That is what makes the quantity portable beyond the benchmarks that happen to contain
replicates.

\section{What this does and does not imply}\label{sec:scope}

It does not indict results obtained by querying these benchmarks in the standard way, since that protocol
reads the published three-run mean rather than a single run. Being imprecise here would be the fastest
route to having the work dismissed.

It bears directly on four things. Methods that train architectures themselves and report the outcome,
where a single run is the evidence. Predictor and proxy evaluation, where the label is noisy and the
ceiling is unreported, so an absolute correlation cannot be interpreted. Multi-fidelity and
early-stopping search, which Section~\ref{sec:results} shows is ranking on a signal that carries no
ordering information at 4 and 12 epochs. And the benchmarks' own reported optima together with every
regret curve measured against them.

\section{Limitations}\label{sec:limits}

The phenomenon is a folk fact and we claim only its quantification. \pending{cite Surrogate NAS
Benchmarks, which models evaluation noise on 500 architectures with 5 seeds and introduces sparse Kendall
tau that rounds accuracies to 0.1 per cent precisely to ignore what we measure; Bouthillier et al.\ 2021
for the general benchmark case; Yang et al.\ 2020 §5.3 as the nearest prior art on ground truth against
itself; Summers and Dinneen 2021. None of these has been read in full and none of the sentences citing
them may stand until they are.}

The NAS-Bench-201 numbers rest on a converted release that fills never-evaluated architecture--seed cells
with the mean of the remaining seeds. Imputed cells have artificially low variance, which deflates both
the noise estimates and the flip rates by an amount we have not quantified and cannot identify from the
converted file. This is the most attackable fact in the paper. NAS-Bench-101 has no such imputation.

Pairs share architectures, so pair-level rates are not independent observations and the intervals should
come from an architecture-level bootstrap, which is not yet done. Everything is measured on validation
accuracy while papers report test accuracy. The flip rates use uniformly drawn pairs, whereas searches
operate at the top of the distribution, so those rows are owed restricted to it. And a per-architecture
$\sigma_W$ from three runs has two degrees of freedom, so it is good to roughly $\pm 60$ per cent;
statements about individual architectures are correspondingly weak, while pooled statements are not.

Finally, the residual of $-0.12$ in NAS-Bench-101's top-1000 and top-10000 cells is unexplained.
Excluding the collapse-prone architectures does not move it, since good architectures do not collapse.
Non-normality of quality at the top and the two to three per cent of tied pairs there are the candidates.
\end{document}
